\chapter{Low-Dimensional Topology}\label{low-dimensional-topology}
\chapterstatus{planned}

In dimensions two, three and four topology has a flavour of its own: surfaces are classified completely,
knots in three-space carry a rich algebra of invariants, three-manifolds are governed by geometrisation,
and four-manifolds behave differently in the smooth and in the topological category. This chapter
develops the first two subjects with proofs and surveys the last two.

\noindent\textbf{Prerequisites.} Chapter~\ref{general-topology} (General Topology), Chapter~\ref{homotopy} (Homotopy and Covering Spaces), Chapter~\ref{singular-homology} (Singular Homology and Cohomology), Chapter~\ref{poincare-duality} (Topological Manifolds and Poincar\'e Duality) and Chapter~\ref{groups} (Groups).

\noindent\textbf{Planned contents.}
\begin{enumerate}
\item Surfaces: triangulations, the classification of compact surfaces with and without boundary, Euler
characteristic and orientability; the fundamental group of a surface and its presentation.
\item Curves on surfaces and the mapping class group: Dehn twists, the Dehn--Lickorish theorem, the
mapping class group of the torus is $\mathrm{SL}_2(\mathbb{Z})$; the Nielsen--Thurston classification (survey).
\item Knots and links: tame knots, Reidemeister's theorem, the linking number, tricolourability and Fox
colourings; the knot group and the Wirtinger presentation; the unknot detected by $\pi_1$ (statement).
\item Seifert surfaces and the genus of a knot; the Seifert matrix, the Alexander polynomial and the
Conway skein relation; signature and determinant.
\item The Kauffman bracket and the Jones polynomial; the Tait conjectures on alternating knots (the span
of the Jones polynomial); torus knots and their invariants.
\item Braids: the braid group, Artin's presentation, Alexander's and Markov's theorems.
\item Three-manifolds: Heegaard splittings, lens spaces and their classification up to homotopy and
homeomorphism, Dehn surgery and the Lickorish--Wallace theorem, the Poincar\'e homology sphere.
\item Prime decomposition (Kneser--Milnor), Seifert fibred spaces, Thurston's eight geometries;
geometrisation and the Poincar\'e conjecture (Perelman), statements.
\item Four-manifolds (survey): the intersection form, Freedman's classification, Donaldson's theorem,
exotic $\mathbb{R}^4$.
\item More knot invariants: the HOMFLY-PT polynomial, Vassiliev (finite type) invariants and chord diagrams, the Kontsevich integral (survey); knot tables and the knots up to eight crossings.
\item Concordance and slice knots, the slice genus, the Milnor conjecture on the unknotting number of torus knots (statement); hyperbolic knots, the figure-eight knot complement, hyperbolic volume,
Mostow rigidity (statement), the volume conjecture (survey).
\item Knot homologies and quantum invariants (survey): Khovanov homology, knot Floer homology, Chern--Simons theory and the Witten--Reshetikhin--Turaev invariants of $3$-manifolds.
\item Exercises with solutions.
\end{enumerate}
